\section{Threat Model and Motivation}
\label{sec:threat}

We consider a white-box adversary that knows the target model, the draft
model, the tokenizer, and the deployed speculative verification rule. At
runtime, however, the adversary remains strictly prompt-bounded: the attacker
may modify only a short assistant-boundary prefix in the prompt. The attacker
cannot change model parameters, decoding temperature, token limits, system
prompts, verification logic, scheduler behavior, or hardware placement at
attack time.

The attacker seeks a discrete prefix $z \in \mathcal{V}^m$ that maximizes
speculative inefficiency while preserving user-visible task quality. At the
systems level, the objective is to increase the verifier work required per
request. We operationalize this through counter-based metrics such as aggregate
acceptance and block efficiency together with wall-clock latency. Formally, the
attacker seeks a short prefix that solves a constrained problem of the form
\[
\max_{z \in \mathcal{V}^m} \;\; \mathcal{L}_{\mathrm{cost}}(z)
\qquad
\text{s.t.}
\qquad
\mathcal{Q}(z) \geq \mathcal{Q}_{\mathrm{benign}} - \delta,
\]
where $\mathcal{L}_{\mathrm{cost}}$ denotes speculative inefficiency and
$\mathcal{Q}$ denotes user-visible task quality.

Why does such an attack matter? In multi-tenant LLM serving, speculative
decoding is deployed precisely because it increases throughput per GPU. When an
adversarial prefix collapses verifier acceptance, the system is forced back
toward a more sequential decoding regime. The immediate victim is the attacked
request, but the operational consequence is broader: extra verifier work
consumes shared accelerator time, reduces node throughput, and degrades the
quality of service seen by co-located users. In that sense, prompt-side
acceptance collapse functions as a resource-exhaustion attack on the
speculative serving stack, even though the adversary changes only the input
text and not the serving infrastructure itself.

The remaining challenge is methodological. The attacker needs a discrete prefix
that works in the deployed verifier, but the most powerful search tools are
continuous. The next section addresses that obstacle directly by formalizing
the projection gap and constructing an optimizer that treats it as part of the
attack objective rather than as an afterthought.
